\chapter{Capítulo IV: Diseño}
\section{Diseño General del Sistema}

El sistema \textbf{VP Planillas} fue concebido bajo una arquitectura de software cliente-servidor, estructurada en \textbf{capas lógicas} que separan las responsabilidades funcionales. Esta decisión permite una mejor escalabilidad, mantenibilidad y posibilidad de integración con futuros sistemas.

\subsection{Arquitectura del Sistema}

El sistema sigue una arquitectura en \textit{n capas}, compuesta por:

\begin{itemize}
	\item \textbf{Capa de presentación (Frontend):} desarrollada con \textit{Next.js} y \textit{TypeScript}, proporciona una interfaz amigable y responsiva.
	\item \textbf{Capa de lógica de negocio (Backend):} construida con \textit{Node.js} y \textit{Express.js}. Gestiona la lógica, validaciones y flujos de datos.
	\item \textbf{Capa de datos (Persistencia):} implementada con \textit{PostgreSQL} y gestionada mediante \textit{Prisma ORM}.
\end{itemize}

La comunicación entre frontend y backend se maneja mediante APIs RESTful protegidas con cifrado HTTPS (TLS).

\subsection{Módulos Funcionales}

El sistema contempla los siguientes módulos principales:

\begin{itemize}
	\item Autenticación y control de acceso.
	\item Gestión de empleados (CRUD).
	\item Cálculo automático de nómina.
	\item Registro de eventos laborales (vacaciones, incapacidades, permisos).
	\item Generación de reportes oficiales.
	\item Emisión de comprobantes de pago.
	\item Configuración del sistema.
\end{itemize}

\subsection{Interfaces de Usuario}

Las interfaces siguen principios de \textit{accesibilidad (WCAG)} y \textit{usabilidad}, con diseño responsivo y tipografía \textit{INTER}, utilizando React y Tailwind CSS. Las vistas incluyen:

\begin{itemize}
	\item Inicio de sesión.
	\item Panel principal.
	\item Gestión de empleados.
	\item Cálculo de nómina.
	\item Gestión de eventos laborales.
	\item Reportes.
	\item Comprobantes de pago.
	\item Configuración del sistema.
\end{itemize}

\subsection{Estándares y Convenciones Aplicadas}

\begin{itemize}
	\item \textbf{Base de datos:} tablas con prefijo \texttt{VPG\_}, campos en \texttt{snake\_case}, normalización hasta 3FN.
	\item \textbf{Código fuente:} uso de \texttt{TypeScript}, control de versiones con Git, repositorio en GitHub, estructura por dominios.
	\item \textbf{Seguridad:} cifrado de contraseñas, tokens de autenticación, roles de usuario, bitácora de trazabilidad.
	\item \textbf{Pruebas:} funcionales por sprint, revisiones por pares, validaciones con usuarios reales.
\end{itemize}

\subsection{Diseño de Base de Datos}

El diseño de base de datos del sistema está basado en una estructura relacional robusta que sigue principios de normalización y nomenclatura estandarizada. Las entidades, relaciones, claves primarias y foráneas han sido cuidadosamente definidas para asegurar integridad y eficiencia.

\begin{itemize}
	\item \textbf{Modelo Relacional:} Se encuentra en el Anexo 1 y muestra gráficamente todas las tablas del sistema.
	\item \textbf{Scripts - Backup:} Incluidos en el Anexo 2, con el código SQL para la creación e inserción de datos base.
	\item \textbf{Diccionario de datos:} Documentado en el Anexo 3, describe los atributos, tipos de datos, restricciones y relaciones de cada tabla.
\end{itemize}

\textbf{Nota:} Todos estos elementos técnicos han sido incorporados como anexos para facilitar su consulta y validación por parte del equipo evaluador.


\subsection{Consideraciones de Integración}

\begin{itemize}
	\item Integración con reloj marcador vía archivos CSV/Excel o API.
	\item Emisión de reportes en formatos \texttt{PDF} o \texttt{Excel} mediante bibliotecas como \texttt{PDFKit} o \texttt{jsPDF}.
\end{itemize}
